% Options for packages loaded elsewhere
\PassOptionsToPackage{unicode}{hyperref}
\PassOptionsToPackage{hyphens}{url}
\PassOptionsToPackage{dvipsnames,svgnames,x11names}{xcolor}
\PassOptionsToPackage{space}{xeCJK}
\documentclass[
]{article}
\usepackage{xcolor}
\usepackage[margin=2.5cm]{geometry}
\usepackage{amsmath,amssymb}
\setcounter{secnumdepth}{-\maxdimen} % remove section numbering
\usepackage{iftex}
\ifPDFTeX
  \usepackage[T1]{fontenc}
  \usepackage[utf8]{inputenc}
  \usepackage{textcomp} % provide euro and other symbols
\else % if luatex or xetex
  \usepackage{unicode-math} % this also loads fontspec
  \defaultfontfeatures{Scale=MatchLowercase}
  \defaultfontfeatures[\rmfamily]{Ligatures=TeX,Scale=1}
\fi
\usepackage{lmodern}
\ifPDFTeX\else
  % xetex/luatex font selection
    \setmainfont[]{DejaVu Sans}
    \setmonofont[]{DejaVu Sans Mono}
  \ifXeTeX
    \usepackage{xeCJK}
    \setCJKmainfont[]{Noto Sans CJK SC}
          \fi
  \ifLuaTeX
    \usepackage[]{luatexja-fontspec}
    \setmainjfont[]{Noto Sans CJK SC}
  \fi
\fi
% Use upquote if available, for straight quotes in verbatim environments
\IfFileExists{upquote.sty}{\usepackage{upquote}}{}
\IfFileExists{microtype.sty}{% use microtype if available
  \usepackage[]{microtype}
  \UseMicrotypeSet[protrusion]{basicmath} % disable protrusion for tt fonts
}{}
\makeatletter
\@ifundefined{KOMAClassName}{% if non-KOMA class
  \IfFileExists{parskip.sty}{%
    \usepackage{parskip}
  }{% else
    \setlength{\parindent}{0pt}
    \setlength{\parskip}{6pt plus 2pt minus 1pt}}
}{% if KOMA class
  \KOMAoptions{parskip=half}}
\makeatother
\usepackage{longtable,booktabs,array}
\usepackage{calc} % for calculating minipage widths
% Correct order of tables after \paragraph or \subparagraph
\usepackage{etoolbox}
\makeatletter
\patchcmd\longtable{\par}{\if@noskipsec\mbox{}\fi\par}{}{}
\makeatother
% Allow footnotes in longtable head/foot
\IfFileExists{footnotehyper.sty}{\usepackage{footnotehyper}}{\usepackage{footnote}}
\makesavenoteenv{longtable}
\setlength{\emergencystretch}{3em} % prevent overfull lines
\providecommand{\tightlist}{%
  \setlength{\itemsep}{0pt}\setlength{\parskip}{0pt}}
\usepackage{bookmark}
\IfFileExists{xurl.sty}{\usepackage{xurl}}{} % add URL line breaks if available
\urlstyle{same}
\hypersetup{
  colorlinks=true,
  linkcolor={Maroon},
  filecolor={Maroon},
  citecolor={Blue},
  urlcolor={Blue},
  pdfcreator={LaTeX via pandoc}}

\author{}
\date{}

\begin{document}

\section{基点穿折越理论泛化：如何选择最快收敛的基点，构造能充分生成人类数域空间的基点，并让每个空间都不处于无穷远}\label{ux57faux70b9ux7a7fux6298ux8d8aux7406ux8bbaux6cdbux5316ux5982ux4f55ux9009ux62e9ux6700ux5febux6536ux655bux7684ux57faux70b9ux6784ux9020ux80fdux5145ux5206ux751fux6210ux4ebaux7c7bux6570ux57dfux7a7aux95f4ux7684ux57faux70b9ux5e76ux8ba9ux6bcfux4e2aux7a7aux95f4ux90fdux4e0dux5904ux4e8eux65e0ux7a77ux8fdc}

\textbf{Basepoint-Teleport (Worn-Zhe-Yue) Generalization: Choosing the
Fastest-Converging Basepoint, Designing a Basepoint that Generates the
Human Number-Field Space, and Removing Every Space from Infinity}

\emph{2026-08-12 · Internal research paper · 方法基础: 基点穿折越
(Worn-Zhe-Yue)}

\begin{center}\rule{0.5\linewidth}{0.5pt}\end{center}

\subsection{Abstract}\label{abstract}

We address the basepoint-selection problem in the basepoint-teleport
(穿折越, Worn-Zhe-Yue) framework: given a structure (here: the human
number-field space), how do we choose the basepoint that (a) converges
fastest to the basepoint-invariant structure, (b) is capable of
\emph{generating} the full number-field chain ℕ → ℤ → ℚ → ℝ → ℂ, and (c)
removes every axis so that no space remains at infinity? We show that
the number-field chain is a \emph{convergence hierarchy}: each level
incorporates the projection artifacts of the previous level as explicit
structure (ℕ cannot see subtraction → ℤ adds inverses; ℤ cannot see
division → ℚ adds fractions; ℚ cannot see limits → ℝ
completes/compactifies, drawing ∞ into a finite circle; ℝ cannot see √-1
→ ℂ takes the algebraic closure). The basepoint 1 (the unit) generates
the entire chain through the operations of successor (ℕ), inverse (ℤ),
reciprocal (ℚ), limit (ℝ), and rotation by 90° (ℂ, i = √-1). We define
\emph{axis peeling}: the removal of representation
(projection-direction) dependence, so that every structure is expressed
in basepoint-invariant coordinates. We give the fastest-convergence
criterion: a basepoint converges fastest when it exposes the generating
structure directly (coordinates are integers, or structure constants, or
symmetry-adapted), so that the projection artifacts (irrationals,
transcendentals, asymmetries) are recognized as such rather than carried
as pseudo-properties. Finally, \emph{no space remains at infinity}: each
level is compactified (ℝ → S¹, ℂ → S²), so every space is finite and
closed under its defining operation.

\subsection{1. Introduction}\label{introduction}

\subsubsection{1.1 Motivation}\label{motivation}

The number line is a pointed, directed structure: basepoint 0, unit 1,
direction positive/negative. The basepoint-teleport (穿折越) framework
treats a basepoint as a genuine degree of freedom: choosing a basepoint
stabilizes relative structure while drifting the codomain. This paper
asks the \emph{design} question: which basepoint should we choose, and
how should we design a basepoint that can generate the whole human
number system?

Three sub-questions structure the paper: 1. \textbf{Fastest
convergence}: which basepoint makes the structure converge fastest to
its basepoint-invariant (axis-free) form? 2. \textbf{Generativity}:
which basepoint can generate ℕ → ℤ → ℚ → ℝ → ℂ, the full number-field
space? 3. \textbf{No infinity}: how do we peel away the axes so that no
space sits at infinity?

\subsubsection{1.2 Contributions}\label{contributions}

\begin{enumerate}
\def\labelenumi{\arabic{enumi}.}
\tightlist
\item
  \textbf{The number-field chain as a convergence hierarchy}
  (Observation 2.1): each level ℕ → ℤ → ℚ → ℝ → ℂ incorporates the
  previous level's projection artifacts as explicit structure.
\item
  \textbf{The basepoint 1 generates the full chain} (Proposition 3.1):
  through successor (ℕ), inverse (ℤ), reciprocal (ℚ), limit (ℝ),
  rotation 90° (ℂ). The unit is the universal generator; i is the y-axis
  choice that closes the algebra.
\item
  \textbf{Axis peeling} (Definition 4.1): removing projection-direction
  dependence; a structure's axis-free form is its basepoint-invariant
  (D-space) form.
\item
  \textbf{Fastest-convergence criterion} (Principle 4.2): choose the
  basepoint exposing the generating structure directly, so projection
  artifacts are recognized rather than carried.
\item
  \textbf{Compactification removes infinity} (Observation 5.1): every
  level is compactified (ℝ → S¹, ℂ → S²); ``not at infinity'' means
  drawing ∞ into the finite structure.
\end{enumerate}

\subsubsection{1.3 Honest boundary}\label{honest-boundary}

The individual facts are classical (KNOWN): the number-field tower,
one-point compactification, algebraic closure. The contribution is the
\emph{unified basepoint-teleport design account}: the number system as a
convergence hierarchy under basepoint choice, the unit as universal
generator, the fastest-convergence criterion, and axis peeling. This
unified account is original to this framework (NO\_PRIOR\_RESULT\_FOUND
in the searched sources).

\subsection{2. The number-field chain as a convergence
hierarchy}\label{the-number-field-chain-as-a-convergence-hierarchy}

\subsubsection{2.1 The levels and their blind
spots}\label{the-levels-and-their-blind-spots}

Each number system ℕ, ℤ, ℚ, ℝ, ℂ has a \emph{blind spot}: an operation
that is not closed within it. The next level is precisely the
incorporation of that operation.

\textbf{Observation 2.1} (the convergence hierarchy). The chain ℕ → ℤ →
ℚ → ℝ → ℂ is a hierarchy in which each level incorporates the previous
level's blind spot as explicit structure:

\begin{longtable}[]{@{}
  >{\raggedright\arraybackslash}p{(\linewidth - 4\tabcolsep) * \real{0.3333}}
  >{\raggedright\arraybackslash}p{(\linewidth - 4\tabcolsep) * \real{0.3333}}
  >{\raggedright\arraybackslash}p{(\linewidth - 4\tabcolsep) * \real{0.3333}}@{}}
\toprule\noalign{}
\begin{minipage}[b]{\linewidth}\raggedright
Level
\end{minipage} & \begin{minipage}[b]{\linewidth}\raggedright
Blind spot (projection artifact)
\end{minipage} & \begin{minipage}[b]{\linewidth}\raggedright
Incorporated by next level
\end{minipage} \\
\midrule\noalign{}
\endhead
\bottomrule\noalign{}
\endlastfoot
ℕ & cannot see subtraction & ℤ adds additive inverses \\
ℤ & cannot see division & ℚ adds multiplicative inverses (fractions) \\
ℚ & cannot see limits & ℝ completes (Cauchy limits), draws ∞ into S¹ \\
ℝ & cannot see √-1 & ℂ takes algebraic closure, draws ∞ into S² \\
\end{longtable}

\textbf{Observation 2.2} (projection reading). The blind spots are
\emph{projection artifacts} of the lower representation: ``subtraction
doesn't close in ℕ'' is the artifact of the ℕ-axis (no negative
direction); ``division doesn't close in ℤ'' is the artifact of the
ℤ-axis (no reciprocal); ``limits don't close in ℚ'' is the artifact of
the ℚ-axis (no completion); ``√-1 doesn't exist in ℝ'' is the artifact
of the ℝ-axis (no perpendicular direction). Each level peels away one
axis artifact and becomes more closed, more basepoint-independent.

\subsection{3. Designing the generating
basepoint}\label{designing-the-generating-basepoint}

\subsubsection{3.1 The unit 1 as universal
generator}\label{the-unit-1-as-universal-generator}

\textbf{Proposition 3.1} (the unit generates the full chain). Starting
from the basepoint 1 and the unit element, the full number-field chain
is generated by successive operations:

\begin{itemize}
\tightlist
\item
  \textbf{ℕ}: successor (add 1 repeatedly) --- 1, 2, 3, \ldots{}
\item
  \textbf{ℤ}: inverse (subtract 1, negate) --- 0, -1, -2, \ldots{}
\item
  \textbf{ℚ}: reciprocal (1/n) --- 1/2, 1/3, \ldots{}
\item
  \textbf{ℝ}: limit (Cauchy completion, Dedekind) --- √2 as the limit of
  1, 1.4, 1.41, 1.414, \ldots{}
\item
  \textbf{ℂ}: rotation by 90° (i = √-1, i² = -1) --- the algebraic
  closure
\end{itemize}

\emph{The unit 1 is the universal basepoint: every number in every field
is generated from it by the above operations.}

\subsubsection{3.2 The 90° axis as the algebraic-closure
basepoint}\label{the-90-axis-as-the-algebraic-closure-basepoint}

\textbf{Observation 3.2} (i is a y-axis choice). The imaginary unit i is
the unit vector of the y-axis (90° rotation). Choosing the y-axis is
choosing whether the system is algebraically closed: without the 90°
direction (the y-axis), ℝ cannot see √-1; with it, ℂ closes.

\textbf{Observation 3.3} (y-axis = basepoint for closure). The y-axis is
a basepoint in the C007 sense: it determines the codomain (whether the
field is closed under √-1), while the underlying algebraic structure is
basepoint-independent. Choosing the y-axis is choosing which projection
of the number system to observe.

\subsubsection{3.3 What a basepoint must carry to generate the number
space}\label{what-a-basepoint-must-carry-to-generate-the-number-space}

\textbf{Design principle 3.4} (a generating basepoint). To generate the
full human number-field space, a basepoint must carry: 1. \textbf{A
unit} (the 1, the seed of successor and of every multiplication); 2.
\textbf{An inverse operation} (yielding negatives and reciprocals); 3.
\textbf{A completion/limit} (yielding reals, drawing ∞ into finite
form); 4. \textbf{A perpendicular direction} (yielding √-1, closing the
algebra).

The basepoint 1 with the operations of
successor/inverse/reciprocal/limit/rotation satisfies all four. This is
the design of a basepoint that \emph{sufficiently constructs the human
number space}.

\subsection{4. Choosing the fastest-converging
basepoint}\label{choosing-the-fastest-converging-basepoint}

\subsubsection{4.1 Convergence defined}\label{convergence-defined}

\textbf{Definition 4.1} (axis peeling). \emph{Axis peeling} is the
removal of representation (projection-direction) dependence: expressing
a structure in coordinates that depend only on the basepoint-invariant
structure (the D-space of C007), not on the choice of axis direction.

\textbf{Definition 4.2} (convergence). A structure \emph{converges} when
axis peeling reaches a fixed point: the coordinates no longer change
when the axis (projection direction) is further varied --- i.e., the
codomain equals the basepoint-invariant structure, with no remaining
projection artifacts.

\subsubsection{4.2 The fastest-convergence
criterion}\label{the-fastest-convergence-criterion}

\textbf{Principle 4.3} (fastest convergence). A basepoint converges
fastest when it \emph{exposes the generating structure directly}: the
coordinates are the integers (for axes), or the structure constants (for
circles), or symmetry-adapted (for groups). In such coordinates,
projection artifacts (irrationals, transcendentals, asymmetries) are
recognized \emph{as} artifacts rather than carried as pseudo-properties.

\begin{itemize}
\tightlist
\item
  \textbf{θ-axis}: choose the axis itself as the x'-axis ⟹ coordinates
  are integers; the ``irrational'' n·cos θ is recognized as a projection
  (fastest convergence: one rotation).
\item
  \textbf{Real line}: choose the circle (S¹) ⟹ π, e are structure
  constants; their transcendence is recognized as a line-representation
  artifact (fastest convergence: one compactification).
\item
  \textbf{NS}: choose O(n) over SO(n) ⟹ time reversal T is a group
  element; the ``irreversibility'' is recognized as an
  SO(n)-representation artifact (fastest convergence: one symmetry
  enlargement).
\end{itemize}

\textbf{Observation 4.4} (slow convergence). Choosing a basepoint that
does not expose the generating structure carries artifacts as
pseudo-properties: the real line ``observes'' √2 and π as if they were
intrinsic irrationals; a basepoint adapted to the generating structure
reveals them as projections. The former converges slowly (each artifact
must be peeled one at a time); the latter converges in one step.

\subsubsection{4.3 The convergence hierarchy as peeling
order}\label{the-convergence-hierarchy-as-peeling-order}

\textbf{Observation 4.5} (the number chain is the peeling sequence). The
chain ℕ → ℤ → ℚ → ℝ → ℂ is exactly the axis-peeling sequence: each step
removes one axis artifact (subtraction, division, limits, √-1), and the
fastest-converging basepoint is the one at the top of the chain (ℂ, the
algebraic closure, where all blind spots are incorporated and the space
is closed). But the \emph{design} of that basepoint requires passing
through all levels --- the convergence is fastest precisely when every
level's artifact has been incorporated.

\subsection{5. No space remains at
infinity}\label{no-space-remains-at-infinity}

\subsubsection{5.1 Compactification draws ∞ into the finite
structure}\label{compactification-draws-into-the-finite-structure}

\textbf{Observation 5.1} (every space is compactified). Each level of
the number chain is compactified, so no space sits at infinity:

\begin{itemize}
\tightlist
\item
  \textbf{ℝ}: one-point compactification ≃ S¹ (the real line closes into
  a circle; +∞ and -∞ meet at one point).
\item
  \textbf{ℂ}: one-point compactification ≃ S² (the Riemann sphere; ∞ is
  the north pole).
\item
  \textbf{The circle axis (S¹)}: on it, the structure constants π and e
  are finite, intrinsic, and no number is ``at infinity.''
\end{itemize}

\textbf{Definition 5.2} (not at infinity). A space is \emph{not at
infinity} when it is compactified: every divergent point is drawn into
the finite structure as a boundary or a meeting point, so no coordinate
can escape the space.

\textbf{Observation 5.3} (peeling removes infinity). Axis peeling and
compactification coincide: peeling the line's axis artifacts
(irrationality of √2, transcendence of π, asymmetry of time) is the same
as drawing those artifacts into explicit finite structure (projections
along an axis, constants of a circle, group elements). A fully peeled
space is finite, closed, and contains no ∞.

\subsection{6. The design recipe: a basepoint for the human number
space}\label{the-design-recipe-a-basepoint-for-the-human-number-space}

\textbf{Design 6.1} (the recipe). To construct a basepoint that
generates the full human number space and makes every space converge
(finite, not at infinity):

\begin{enumerate}
\def\labelenumi{\arabic{enumi}.}
\tightlist
\item
  \textbf{Start from the unit 1} --- the seed of successor, inverse,
  reciprocal, limit, rotation.
\item
  \textbf{Add the inverse operation} (peel subtraction/division): ℤ, ℚ.
\item
  \textbf{Complete} (peel limits): ℝ; compactify to S¹ (draw ∞ in).
\item
  \textbf{Take the 90° direction} (peel √-1): ℂ; compactify to S² (draw
  ∞ in).
\item
  \textbf{Choose the fastest-converging coordinates} at each level: the
  generating-structure-adapted axis (integers along the θ-axis;
  structure constants on the circle; symmetry-adapted for groups).
\end{enumerate}

The resulting space: ℂ, algebraically closed, compact (Riemann sphere),
with every number expressible as a finite combination of basepoint-1
operations --- no irrationality, no transcendence, no asymmetry, no
infinity, except as recognized projection artifacts of a chosen
(non-adapted) axis.

\subsection{7. Conclusion}\label{conclusion}

The human number-field space ℕ → ℤ → ℚ → ℝ → ℂ is a convergence
hierarchy under basepoint choice: each level incorporates the previous
level's projection artifacts (subtraction, division, limits, √-1) as
explicit structure. The basepoint 1, together with the operations of
successor/inverse/reciprocal/limit/90°-rotation, generates the full
chain. The fastest-converging basepoint is the one that exposes the
generating structure directly, so that irrationals, transcendentals, and
asymmetries are recognized as projection artifacts rather than carried
as pseudo-properties. Axis peeling removes representation dependence;
compactification (ℝ → S¹, ℂ → S²) draws ∞ into the finite structure. The
fully peeled space is finite, closed, and contains no point at infinity
--- a basepoint-invariant, axis-free, convergent design of the number
system.

\subsection{Appendix A: identity checks
(恒等式先验)}\label{appendix-a-identity-checks-ux6052ux7b49ux5f0fux5148ux9a8c}

\begin{itemize}
\tightlist
\item
  2·cos 45° = √2 (verified: 1.414214); 2·cos 30° = √3 (verified:
  1.732051)
\item
  e\^{}(i·2πk) = 1 for all k ∈ ℤ (verified); e\^{}(i·(2k+1)π) = -1 for
  all k ∈ ℤ (verified)
\item
  Euler: e\^{}(iπ) = -1 (verified: -1.000000 + 0.000000i)
\item
  i² = -1 (algebraic closure of ℝ); Cauchy sequence 1, 1.4, 1.41, 1.414,
  \ldots{} → √2
\item
  Compactifications: ℝ ∪ \{∞\} ≃ S¹; ℂ ∪ \{∞\} ≃ S² (classical, KNOWN)
\end{itemize}

\subsection{Appendix B: honesty
boundaries}\label{appendix-b-honesty-boundaries}

\begin{itemize}
\tightlist
\item
  All individual mathematical facts are classical (KNOWN): number-field
  tower, algebraic closure, one-point compactification, Euler's
  formula/identity.
\item
  The \emph{unified basepoint-teleport design account} --- the number
  system as a convergence hierarchy, the unit as universal generator,
  the fastest-convergence criterion, axis peeling, and ``no space at
  infinity'' via compactification --- is the paper's contribution
  (NO\_PRIOR\_RESULT\_FOUND in the searched sources).
\item
  This is a conceptual framework paper, not a theorem-proving paper
  (Lean formalization of the axis-peeling construction is future work).
\end{itemize}

\end{document}
